% Source: https://artofproblemsolving.com/community/c3h1635553_sierpinski_triangle
% Stable identity: aops-topic-1635553-q01
The Sierpinski triangle is built in steps. Start (step $0$) with a shaded equilateral triangle of side $1$. At each step, every shaded triangle is divided into four equal equilateral triangles by joining the midpoints of its sides, and the middle one is removed (unshaded).
\begin{enumerate}
\item[(a)] How many shaded triangles are there after step $n$, and what is the side length of each?
\item[(b)] What fraction of the original area is still shaded after step $n$?
\item[(c)] What is the total perimeter of the shaded triangles after step $n$?
\item[(d)] Show that the shaded area can be made smaller than any given positive number, while the total perimeter becomes larger than any given number.
\end{enumerate}
